%********************************************%
%*       Generated from PreTeXt source      *%
%*       on 2026-08-31T19:30:26Z       *%
%*   A recent stable commit (2022-07-01):   *%
%* 6c761d3dba23af92cba35001c852aac04ae99a5f *%
%*                                          *%
%*         https://pretextbook.org          *%
%*                                          *%
%********************************************%
\documentclass[twoside,14pt,]{extbook}
%% Custom Preamble Entries, early (use latex.preamble.early)
%% Always open on odd page
%%   The following adjusts cleardoublepage to remove twosided
%%   check so that we open on odd pages even in one-sided mode
%%   by adding an extra blank page on the preceding even page.
\makeatletter%
\def\cleardoublepage{%
\clearpage\ifodd\c@page\else\thispagestyle{empty}\hbox{}\newpage\if@twocolumn\hbox{}\newpage\fi\fi%
}
\makeatother%
%% Default LaTeX packages
%%   1.  always employed (or nearly so) for some purpose, or
%%   2.  a stylewriter may assume their presence
\usepackage{geometry}
%% Some aspects of the preamble are conditional,
%% the LaTeX engine is one such determinant
\usepackage{ifthen}
%% etoolbox has a variety of modern conveniences
\usepackage{etoolbox}
\usepackage{ifxetex,ifluatex}
%% Raster graphics inclusion
\usepackage{graphicx}
%% Color support, xcolor package
%% Always loaded, for: add/delete text, author tools
%% Here, since tcolorbox loads tikz, and tikz loads xcolor
\PassOptionsToPackage{dvipsnames,svgnames,table}{xcolor}
\usepackage{xcolor}
%% begin: defined colors, via xcolor package, for styling
%% end: defined colors, via xcolor package, for styling
%% Colored boxes, and much more, though mostly styling
%% skins library provides "enhanced" skin, employing tikzpicture
%% boxes may be configured as "breakable" or "unbreakable"
%% "raster" controls grids of boxes, aka side-by-side
\usepackage{tcolorbox}
\tcbuselibrary{skins}
\tcbuselibrary{breakable}
\tcbuselibrary{raster}
%% We load some "stock" tcolorbox styles that we use a lot
%% Placement here is provisional, there will be some color work also
%% First, black on white, no border, transparent, but no assumption about titles
\tcbset{ bwminimalstyle/.style={size=minimal, boxrule=-0.3pt, frame empty,
colback=white, colbacktitle=white, coltitle=black, opacityfill=0.0} }
%% Second, bold title, run-in to text/paragraph/heading
%% Space afterwards will be controlled by environment,
%% independent of constructions of the tcb title
%% Places \blocktitlefont onto many block titles
\tcbset{ runintitlestyle/.style={fonttitle=\blocktitlefont\upshape\bfseries, attach title to upper} }
%% Spacing prior to each exercise, anywhere
\tcbset{ exercisespacingstyle/.style={before skip={1.5ex plus 0.5ex}} }
%% Spacing prior to each block
\tcbset{ blockspacingstyle/.style={before skip={2.0ex plus 0.5ex}} }
%% xparse allows the construction of more robust commands,
%% this is a necessity for isolating styling and behavior
%% The tcolorbox library of the same name loads the base library
\tcbuselibrary{xparse}
%% The tcolorbox library loads TikZ, its calc package is generally useful,
%% and is necessary for some smaller documents that use partial tcolor boxes
%% See:  https://github.com/PreTeXtBook/pretext/issues/1624
\usetikzlibrary{calc}
%% We use some more exotic tcolorbox keys to restore indentation to parboxes
\tcbuselibrary{hooks}
%% Save default paragraph indentation and parskip for use later, when adjusting parboxes
\newlength{\ptxnormalparindent}
\newlength{\ptxnormalparskip}
%% Prescribe the paragraph indent to the value shared with the
%% XSL-FO conversion, so the two agree, rather than accepting the
%% document class default; "\ptxnormalparindent" then captures it
\AtBeginDocument{\setlength{\parindent}{1.5em}}
\AtBeginDocument{\setlength{\ptxnormalparindent}{\parindent}}
\AtBeginDocument{\setlength{\ptxnormalparskip}{\parskip}}
\newcommand{\ptxsetparstyle}{\setlength{\parindent}{\ptxnormalparindent}\setlength{\parskip}{\ptxnormalparskip}}%% An appendix is numbered by a letter, and a bare letter beside a
%% title is ambiguous ("A Trip Abroad"), so a letter composed with
%% a title carries a period ("A. Trip Abroad").  These macros are
%% seams in the division styles, empty until the appendices arrive
%% (chapter-level for a book, section-level for an article, so a
%% multi-part number like "E.1" never gains a period); see the
%% uses at  \appendix
\newcommand{\ptxappendixperiodchapter}{}
\newcommand{\ptxappendixperiodsection}{}
%% Hyperref should be here, but likes to be loaded late
%%
%% Inline math delimiters, \(, \), need to be robust
%% 2016-01-31:  latexrelease.sty  supersedes  fixltx2e.sty
%% If  latexrelease.sty  exists, bugfix is in kernel
%% If not, bugfix is in  fixltx2e.sty
%% See:  https://tug.org/TUGboat/tb36-3/tb114ltnews22.pdf
%% and read "Fewer fragile commands" in distribution's  latexchanges.pdf
\IfFileExists{latexrelease.sty}{}{\usepackage{fixltx2e}}
%% Text height identically 9 inches, text width varies on point size
%% See Bringhurst 2.1.1 on measure for recommendations
%% 75 characters per line (count spaces, punctuation) is target
%% which is the upper limit of Bringhurst's recommendations
\geometry{letterpaper,total={476pt,9.0in}}
%% Custom Page Layout Adjustments (use publisher page-geometry entry)
%% This LaTeX file may be compiled with pdflatex, xelatex, or lualatex executables
%% LuaTeX is not explicitly supported, but we do accept additions from knowledgeable users
%% The conditional below provides  pdflatex  specific configuration last
%% begin: engine-specific capabilities
\ifthenelse{\boolean{xetex} \or \boolean{luatex}}{%
%% begin: xelatex and lualatex-specific default configuration
\ifxetex\usepackage{xltxtra}\fi
%% realscripts is the only part of xltxtra relevant to lualatex 
\ifluatex\usepackage{realscripts}\fi
%% end:   xelatex and lualatex-specific default configuration
}{
%% begin: pdflatex-specific default configuration
%% We assume a PreTeXt XML source file may have Unicode characters
%% and so we ask LaTeX to parse a UTF-8 encoded file
%% This may work well for accented characters in Western language,
%% but not with Greek, Asian languages, etc.
%% When this is not good enough, switch to the  xelatex  engine
%% where Unicode is better supported (encouraged, even)
\usepackage[utf8]{inputenc}
%% end: pdflatex-specific default configuration
}
%% end:   engine-specific capabilities
%%
%% Fonts.  Conditional on LaTex engine employed.
%% Default Text Font: The Latin Modern fonts are
%% "enhanced versions of the [original TeX] Computer Modern fonts."
%% We use them as the default text font for PreTeXt output.
%% Automatic Font Control
%% Portions of a document, are, or may, be affected by defined commands
%% These are perhaps more flexible when using  xelatex  rather than  pdflatex
%% The following definitions are meant to be re-defined in a style, using \renewcommand
%% They are scoped when employed (in a TeX group), and so should not be defined with an argument
\newcommand{\divisionfont}{\relax}
\newcommand{\blocktitlefont}{\relax}
\newcommand{\contentsfont}{\relax}
\newcommand{\pagefont}{\relax}
\newcommand{\tabularfont}{\relax}
\newcommand{\xreffont}{\relax}
\newcommand{\titlepagefont}{\relax}
%%
\ifthenelse{\boolean{xetex} \or \boolean{luatex}}{%
%% begin: font setup and configuration for use with xelatex
%% Generally, xelatex is necessary for non-Western fonts
%% fontspec package provides extensive control of system fonts,
%% meaning *.otf (OpenType), and apparently *.ttf (TrueType)
%% that live *outside* your TeX/MF tree, and are controlled by your *system*
%% (it is possible that a TeX distribution will place fonts in a system location)
%%
%% The fontspec package is the best vehicle for using different fonts in  xelatex
%% So we load it always, no matter what a publisher or style might want
%%
\usepackage{fontspec}
%%
%% begin: xelatex main font ("font-xelatex-main" template)
%% Latin Modern Roman is the default font for xelatex and so is loaded with a TU encoding
%% *in the format* so we can't touch it, only perhaps adjust it later
%% in one of two ways (then known by NFSS names such as "lmr")
%% (1) via NFSS with font family names such as "lmr" and "lmss"
%% (2) via fontspec with commands like \setmainfont{Latin Modern Roman}
%% The latter requires the font to be known at the system-level by its font name,
%% but will give access to OTF font features through optional arguments
%% https://tex.stackexchange.com/questions/470008/
%% where-and-how-does-fontspec-sty-specify-the-default-font-latin-modern-roman
%% http://tex.stackexchange.com/questions/115321
%% /how-to-optimize-latin-modern-font-with-xelatex
%%
%% end:   xelatex main font ("font-xelatex-main" template)
%% begin: xelatex math font ("font-xelatex-math" template)
%% end:   xelatex math font ("font-xelatex-math" template)
%% begin: xelatex mono font ("font-xelatex-mono" template)
%% (conditional on non-trivial uses being present in source)
%% end:   xelatex mono font ("font-xelatex-mono" template)
%% begin: xelatex font adjustments ("font-xelatex-style" template)
%% end:   xelatex font adjustments ("font-xelatex-style" template)
%%
%% Extensive support for other languages
\usepackage{polyglossia}
%% Set main/default language based on pretext/@xml:lang value
%% document language code is "en-US", US English
%% usmax variant has extra hypenation
\setmainlanguage[variant=usmax]{english}
%% Enable secondary languages based on discovery of @xml:lang values
%% Enable fonts/scripts based on discovery of @xml:lang values
%% Western languages should be ably covered by Latin Modern Roman
%% end:   font setup and configuration for use with xelatex
}{%
%% begin: font setup and configuration for use with pdflatex
%% begin: pdflatex main font ("font-pdflatex-main" template)
\usepackage{lmodern}
\usepackage[T1]{fontenc}
%% end:   pdflatex main font ("font-pdflatex-main" template)
%% begin: pdflatex mono font ("font-pdflatex-mono" template)
%% (conditional on non-trivial uses being present in source)
%% end:   pdflatex mono font ("font-pdflatex-mono" template)
%% begin: pdflatex font adjustments ("font-pdflatex-style" template)
%% end:   pdflatex font adjustments ("font-pdflatex-style" template)
%% end:   font setup and configuration for use with pdflatex
}
%% Micromanage spacing, etc.  The named "microtype-options"
%% template may be employed to fine-tune package behavior
\usepackage{microtype}
%% Symbols, align environment, commutative diagrams, bracket-matrix
\usepackage{amsmath}
\usepackage{amscd}
\usepackage{amssymb}
%% allow page breaks within display mathematics anywhere
%% level 4 is maximally permissive
%% this is exactly the opposite of AMSmath package philosophy
%% there are per-display, and per-equation options to control this
%% split, aligned, gathered, and alignedat are not affected
\allowdisplaybreaks[4]
%% allow more columns to a matrix
%% can make this even bigger by overriding with  latex.preamble.late  processing option
\setcounter{MaxMatrixCols}{30}
%%
%% Symbols for use in text mode (i.e. not in math mode)
\usepackage{pifont}
%% Solid geometric shapes to mark ends of environments
\newcommand{\ptxsquare}{\ding{110}}
\newcommand{\ptxtriangle}{\ding{115}}
\newcommand{\ptxdiamond}{\ding{117}}
%% Markers for unordered list items
\newcommand{\ptxlistdisc}{\raisebox{0.3ex}{\textbullet}}
\newcommand{\ptxlistcircle}{\raisebox{0.3ex}{\textopenbullet}}
\newcommand{\ptxlistsquare}{\raisebox{0.4ex}{\scalebox{0.5}{\ding{110}}}}
%%
%% Division Titles, and Page Headers/Footers
%% titlesec package, loading "titleps" package cooperatively
%% See code comments about the necessity and purpose of "explicit" option.
%% The "newparttoc" option causes a consistent entry for parts in the ToC 
%% file, but it is only effective if there is a \titleformat for \part.
%% "pagestyles" loads the  titleps  package cooperatively.
\usepackage[explicit, newparttoc, pagestyles]{titlesec}
%% The companion titletoc package for the ToC.
\usepackage{titletoc}
%% Fixes a bug with transition from chapters to appendices in a "book"
%% See generating XSL code for more details about necessity
\newtitlemark{\chaptertitlename}
%% begin: customizations of page styles via the modal "titleps-style" template
%% Designed to use commands from the LaTeX "titleps" package
%% Plain pages should have the same font for page numbers
\renewpagestyle{plain}{%
\setfoot{}{\pagefont\thepage}{}%
}%
%% Two-page spread as in default LaTeX
\renewpagestyle{headings}{%
\sethead%
[\pagefont\thepage]%
[]
[\pagefont\slshape\MakeUppercase{\ifthechapter{\chaptertitlename\space\thechapter.\space}{}\chaptertitle}]%
{\pagefont\slshape\MakeUppercase{\ifthesection{\thesection.\space\sectiontitle}{}}}%
{}%
{\pagefont\thepage}%
}%
\pagestyle{headings}
%% end: customizations of page styles via the modal "titleps-style" template
%%
%% Create globally-available macros to be provided for style writers
%% These are redefined for each occurence of each division
\newcommand{\divisionnameptx}{\relax}%
\newcommand{\titleptx}{\relax}%
\newcommand{\subtitleptx}{\relax}%
\newcommand{\shortitleptx}{\relax}%
\newcommand{\authorsptx}{\relax}%
\newcommand{\epigraphptx}{\relax}%
%% Create environments for possible occurences of each division
%% Environment for a PTX "preface" at the level of a LaTeX "chapter"
\NewDocumentEnvironment{preface}{mmmmmmm}
{%
\renewcommand{\divisionnameptx}{#1}%
\renewcommand{\titleptx}{#2}%
\renewcommand{\subtitleptx}{#3}%
\renewcommand{\shortitleptx}{#4}%
\renewcommand{\authorsptx}{#5}%
\renewcommand{\epigraphptx}{#6}%
\chapter*{#2}%
\addcontentsline{toc}{chapter}{#4}
\label{#7}%
}{}%
%% Environment for a PTX "chapter" at the level of a LaTeX "chapter"
\NewDocumentEnvironment{chapterptx}{mmmmmmm}
{%
\renewcommand{\divisionnameptx}{#1}%
\renewcommand{\titleptx}{#2}%
\renewcommand{\subtitleptx}{#3}%
\renewcommand{\shortitleptx}{#4}%
\renewcommand{\authorsptx}{#5}%
\renewcommand{\epigraphptx}{#6}%
\chapter[{#4}]{#2}%
\label{#7}%
}{}%
%% Environment for a PTX "section" at the level of a LaTeX "section"
\NewDocumentEnvironment{sectionptx}{mmmmmmm}
{%
\renewcommand{\divisionnameptx}{#1}%
\renewcommand{\titleptx}{#2}%
\renewcommand{\subtitleptx}{#3}%
\renewcommand{\shortitleptx}{#4}%
\renewcommand{\authorsptx}{#5}%
\renewcommand{\epigraphptx}{#6}%
\section[{#4}]{#2}%
\label{#7}%
}{}%
%% Environment for a PTX "handout" at the level of a LaTeX "subsection"
\NewDocumentEnvironment{handout-subsection}{mmmmmmm}
{%
\renewcommand{\divisionnameptx}{#1}%
\renewcommand{\titleptx}{#2}%
\renewcommand{\subtitleptx}{#3}%
\renewcommand{\shortitleptx}{#4}%
\renewcommand{\authorsptx}{#5}%
\renewcommand{\epigraphptx}{#6}%
\subsection[{#4}]{#2}%
\label{#7}%
}{}%
%% Environment for a PTX "handout" at the level of a LaTeX "subsection"
\NewDocumentEnvironment{handout-subsection-numberless}{mmmmmmm}
{%
\renewcommand{\divisionnameptx}{#1}%
\renewcommand{\titleptx}{#2}%
\renewcommand{\subtitleptx}{#3}%
\renewcommand{\shortitleptx}{#4}%
\renewcommand{\authorsptx}{#5}%
\renewcommand{\epigraphptx}{#6}%
\subsection*{#2}%
\addcontentsline{toc}{subsection}{#4}
\label{#7}%
}{}%
%% Environment for a PTX "worksheet" at the level of a LaTeX "section"
\NewDocumentEnvironment{worksheet-section}{mmmmmmm}
{%
\renewcommand{\divisionnameptx}{#1}%
\renewcommand{\titleptx}{#2}%
\renewcommand{\subtitleptx}{#3}%
\renewcommand{\shortitleptx}{#4}%
\renewcommand{\authorsptx}{#5}%
\renewcommand{\epigraphptx}{#6}%
\section[{#4}]{#2}%
\label{#7}%
}{}%
%% Environment for a PTX "worksheet" at the level of a LaTeX "section"
\NewDocumentEnvironment{worksheet-section-numberless}{mmmmmmm}
{%
\renewcommand{\divisionnameptx}{#1}%
\renewcommand{\titleptx}{#2}%
\renewcommand{\subtitleptx}{#3}%
\renewcommand{\shortitleptx}{#4}%
\renewcommand{\authorsptx}{#5}%
\renewcommand{\epigraphptx}{#6}%
\section*{#2}%
\addcontentsline{toc}{section}{#4}
\label{#7}%
}{}%
%%
%% Styles for six traditional LaTeX divisions
\titleformat{\part}[display]
{\divisionfont\Huge\bfseries\centering}{\divisionnameptx\space\thepart}{30pt}{\Huge#1}
[{\Large\centering\authorsptx}]
\titleformat{\chapter}[display]
{\divisionfont\huge\bfseries}{\divisionnameptx\space\thechapter}{20pt}{\Huge#1}
[{\Large\authorsptx}]
\titleformat{name=\chapter,numberless}[display]
{\divisionfont\huge\bfseries}{}{0pt}{#1}
[{\Large\authorsptx}]
\titlespacing*{\chapter}{0pt}{50pt}{40pt}
\titleformat{\section}[hang]
{\divisionfont\Large\bfseries}{\thesection\ptxappendixperiodsection}{1ex}{#1}
[{\large\authorsptx}]
\titleformat{name=\section,numberless}[block]
{\divisionfont\Large\bfseries}{}{0pt}{#1}
[{\large\authorsptx}]
\titlespacing*{\section}{0pt}{3.5ex plus 1ex minus .2ex}{2.3ex plus .2ex}
\titleformat{\subsection}[hang]
{\divisionfont\large\bfseries}{\thesubsection}{1ex}{#1}
[{\normalsize\authorsptx}]
\titleformat{name=\subsection,numberless}[block]
{\divisionfont\large\bfseries}{}{0pt}{#1}
[{\normalsize\authorsptx}]
\titlespacing*{\subsection}{0pt}{3.25ex plus 1ex minus .2ex}{1.5ex plus .2ex}
\titleformat{\subsubsection}[hang]
{\divisionfont\normalsize\bfseries}{\thesubsubsection}{1em}{#1}
[{\small\authorsptx}]
\titleformat{name=\subsubsection,numberless}[block]
{\divisionfont\normalsize\bfseries}{}{0pt}{#1}
[{\normalsize\authorsptx}]
\titlespacing*{\subsubsection}{0pt}{3.25ex plus 1ex minus .2ex}{1.5ex plus .2ex}
\titleformat{\paragraph}[hang]
{\divisionfont\normalsize\bfseries}{\theparagraph}{1em}{#1}
[{\small\authorsptx}]
\titleformat{name=\paragraph,numberless}[block]
{\divisionfont\normalsize\bfseries}{}{0pt}{#1}
[{\normalsize\authorsptx}]
\titlespacing*{\paragraph}{0pt}{3.25ex plus 1ex minus .2ex}{1.5em}
%%
%% Styles for five traditional LaTeX divisions
\titlecontents{part}%
[0pt]{\contentsmargin{0em}\addvspace{1pc}\contentsfont\bfseries}%
{\Large\thecontentslabel\enspace}{\Large}%
{}%
[\addvspace{.5pc}]%
\titlecontents{chapter}%
[0pt]{\contentsmargin{0em}\addvspace{1pc}\contentsfont\bfseries}%
{\large\thecontentslabel\ptxappendixperiodchapter\enspace}{\large}%
{\hfill\bfseries\thecontentspage}%
[\addvspace{.5pc}]%
\titlecontents{section}%
[3.8em]{\contentsfont}%
{\contentslabel[\thecontentslabel\ptxappendixperiodsection]{2.3em}}%
{\hspace*{-2.3em}}%
{\titlerule*[1pc]{.}\contentspage}%
\dottedcontents{subsection}[6.1em]{\contentsfont}{3.2em}{1pc}%
\dottedcontents{subsubsection}[9.3em]{\contentsfont}{4.3em}{1pc}%
%%
%% Begin: Semantic Macros
%% To preserve meaning in a LaTeX file
%%
%% \ptxmono macro for content of "c", "cd", "tag", etc elements
%% Also used automatically in other constructions
%% Simply an alias for \texttt
%% Always defined, even if there is no need, or if a specific tt font is not loaded
\newcommand{\ptxmono}[1]{\texttt{#1}}
%%
%%
%% Following semantic macros are only defined here if their
%% use is required only in this specific document
%%
%% Style of a title on a list item, for ordered and unordered lists
%% Also "task" of exercise, PROJECT-LIKE, EXAMPLE-LIKE
\newcommand{\ptxlititle}[1]{{\slshape#1}}
%% End: Semantic Macros
%% Divisional exercises (and worksheet) as LaTeX environments
%% Third argument is option for extra workspace in worksheets
%% Hanging indent occupies a 5ex width slot prior to left margin
%% Experimentally this seems just barely sufficient for a bold "888."
%% Division exercises, not in exercise group
\tcbset{ divisionexercisestyle/.style={bwminimalstyle, runintitlestyle, exercisespacingstyle, left=5ex, breakable, before upper app={\ptxsetparstyle} } }
\newtcolorbox{divisionexercise}[4]{divisionexercisestyle, before title={\hspace*{-5ex}\makebox[5ex][l]{#1.}}, title={\notblank{#2}{#2}{}}, after title={\notblank{#2}{\space}{}}, phantom={\label{#4}\hypertarget{#4}{}}, after={\notblank{#3}{\par\rule{\ptxworkspacestrutwidth}{#3}\par\vfill}{\par}}}
%% Worksheets and handouts children may have workspaces
\newlength{\ptxworkspacestrutwidth}
%% @workspace strut is invisible
\setlength{\ptxworkspacestrutwidth}{0pt}
%% "tcolorbox" environment for a single image, occupying entire \linewidth
%% arguments are left-margin, width, right-margin, as multiples of
%% \linewidth, and are guaranteed to be positive and sum to 1.0
\tcbset{ imagestyle/.style={bwminimalstyle} }
\NewTColorBox{tcbimage}{mmm}{imagestyle,left skip=#1\linewidth,width=#2\linewidth}
%% Wrapper environment for tcbimage environment with a fourth argument
%% Fourth argument, if nonempty, is a vertical space adjustment
%% and implies image will be preceded by \leavevmode\nopagebreak
%% Intended use is for alignment with a list marker
\NewDocumentEnvironment{image}{mmmm}{\notblank{#4}{\leavevmode\nopagebreak\vspace{#4}}{}\begin{tcbimage}{#1}{#2}{#3}}{\end{tcbimage}%
}
%% More flexible list management, esp. for references
%% But also for specifying labels (i.e. custom order) on nested lists
\usepackage{enumitem}
%% hyperref driver does not need to be specified, it will be detected
\usepackage{hyperref}
%% For a print PDF, no surrounding boxes, so simply textcolor (but still active to preserve spacing)
\hypersetup{hidelinks}
%% A print PDF needs no internal anchors, nor links to them:
%% both dissolve to their content, which preserves spacing
\renewcommand{\hypertarget}[2]{#2}
\renewcommand{\hyperlink}[2]{#2}
%% Less-clever names for hyperlinks are more reliable, *especially* for structural parts
%% See comments in the code to learn more about the importance of this setting
\hypersetup{hypertexnames=false}
%%The  hypertexnames  setting then confuses the hyperlinking from the index
%%This patch resolves the incorrect links, see code for StackExchange post.
\makeatletter
\patchcmd\Hy@EveryPageBoxHook{\Hy@EveryPageAnchor}{\Hy@hypertexnamestrue\Hy@EveryPageAnchor}{}{\fail}
\makeatother
\hypersetup{pdftitle={MATH 1070 - Fall 2026}}
%% If you manually remove hyperref, leave in this next command
%% This will allow LaTeX compilation, employing this no-op command
\providecommand\phantomsection{}
%% Division Numbering: Chapters, Sections, Subsections, etc
%% Division numbers may be turned off at some level ("depth")
%% A section *always* has depth 1, contrary to us counting from the document root
%% The latex default is 3.  If a larger number is present here, then
%% removing this command may make some cross-references ambiguous
%% The precursor variable $numbering-maxlevel is checked for consistency in the common XSL file
\setcounter{secnumdepth}{-1}
%%
%%
%% A faux tcolorbox whose only purpose is to provide common numbering
%% facilities for most blocks (possibly not projects, 2D displays)
%% Controlled by  numbering.theorems.level  processing parameter
\newtcolorbox[auto counter]{block}{}
%% A faux tcolorbox whose only purpose is to provide common numbering
%% facilities for 2D displays which are subnumbered as part of a "sidebyside"
\makeatletter
\newtcolorbox[auto counter, number within=tcb@cnt@block, number freestyle={\noexpand\thetcb@cnt@block(\noexpand\alph{\tcbcounter})}]{subdisplay}{}
\makeatother
%%
%% tcolorbox, with styles, for DEFINITION-LIKE
%%
%% definition: fairly simple numbered block/structure
\tcbset{ definitionstyle/.style={bwminimalstyle, runintitlestyle, blockspacingstyle, after title={\space}, after upper={\space\space\hspace*{\stretch{1}}\ptxdiamond}, before upper app={\ptxsetparstyle}, } }
\newtcolorbox[use counter from=block]{definition}[3]{title={{#1~\thetcbcounter\notblank{#2}{\space\space#2}{}}}, phantomlabel={#3}, breakable, after={\par}, definitionstyle, }
%%
%% tcolorbox, with styles, for EXAMPLE-LIKE
%%
%% example: fairly simple numbered block/structure
\tcbset{ examplestyle/.style={bwminimalstyle, runintitlestyle, blockspacingstyle, after title={\space}, after upper={\space\space\hspace*{\stretch{1}}\ptxtriangle}, before upper app={\ptxsetparstyle}, } }
\newtcolorbox[use counter from=block]{example}[3]{title={{#1~\thetcbcounter\notblank{#2}{\space\space#2}{}}}, phantomlabel={#3}, breakable, after={\par}, examplestyle, }
%%
%% tcolorbox, with styles, for GOAL-LIKE
%%
%% objectives: early in a subdivision, introduction/list/conclusion
\tcbset{ objectivesstyle/.style={bwminimalstyle, blockspacingstyle, fonttitle=\blocktitlefont\large\bfseries, toprule=0.1ex, toptitle=0.5ex, top=2ex, bottom=0.5ex, bottomrule=0.1ex} }
\newtcolorbox{objectives}[2]{title={#1}, phantomlabel={#2}, breakable, objectivesstyle, before upper app={\ptxsetparstyle}}
%%
%% xparse environments for introductions and conclusions of divisions
%%
%% introduction: in a structured division
\NewDocumentEnvironment{introduction}{m}
{\notblank{#1}{\noindent\textbf{#1}\space}{}}{\par\medskip}
%%
%% tcolorbox, with styles, for miscellaneous environments
%%
%% paragraphs: the terminal, pseudo-division
%% We use the lowest LaTeX traditional division
\titleformat{\subparagraph}[runin]{\normalfont\normalsize\bfseries}{\thesubparagraph}{1em}{#1}
\titlespacing*{\subparagraph}{0pt}{3.25ex plus 1ex minus .2ex}{1em}
\NewDocumentEnvironment{paragraphs}{mm}
{\subparagraph*{#1}\label{#2}\hypertarget{#2}{}}{}
%% Graphics Preamble Entries
\usepackage{tikz, pgfplots}
\usetikzlibrary{positioning,matrix,arrows}
\usetikzlibrary{shapes,decorations,shadows,fadings,patterns}
\usetikzlibrary{decorations.markings}
%% If tikz has been loaded, replace ampersand with \amp macro
%% Custom Preamble Entries, late (use latex.preamble.late)
%% Begin: Author-provided macros
%% (From  docinfo/macros  element)
%% Plus three from PTX for XML characters
\newcommand{\N}{\mathbb N}
\newcommand{\Z}{\mathbb Z}
\newcommand{\Q}{\mathbb Q}
\newcommand{\R}{\mathbb R}
\newcommand{\lt}{<}
\newcommand{\gt}{>}
\newcommand{\amp}{&}
%% End: Author-provided macros
\begin{document}
%% bottom alignment is explicit, since it normally depends on oneside, twoside
\raggedbottom
\frontmatter
%% begin: half-title
\thispagestyle{empty}
{\titlepagefont\centering
\vspace*{0.28\textheight}
{\Huge MATH 1070 - Fall 2026}\\[2\baselineskip]
{\LARGE College Algebra}\\
}
\clearpage
%% end:   half-title
%% begin: adcard (empty)
\thispagestyle{empty}
\null%
\clearpage
%% end:   adcard (empty)
%% begin: title page
%% Inspired by Peter Wilson's "titleDB" in "titlepages" CTAN package
\thispagestyle{empty}
{\titlepagefont\centering
\vspace*{0.14\textheight}
%% Target for xref to top-level element is ToC
\addtocontents{toc}{\protect\label{root-1-2}\protect\hypertarget{root-1-2}{}}
{\Huge MATH 1070 - Fall 2026}\\[\baselineskip]
{\LARGE College Algebra}\\[3\baselineskip]
{\Large Matthew Charnley}\\[0.5\baselineskip]
{\Large Wayne State University}\\[3\baselineskip]
{\Large Last Updated: August 31, 2026}\\}
\clearpage
%% end:   title page
%% begin: copyright-page
\thispagestyle{empty}
\vspace*{\stretch{2}}
\vspace*{\stretch{1}}
\null\clearpage
%% end:   copyright-page
%
%
\typeout{************************************************}
\typeout{Preface  Preface}
\typeout{************************************************}
%
\begin{preface}{Preface}{Preface}{}{Preface}{}{}{preface}
\begin{paragraphs}{Introduction.}{preface-1}%
This eBook contains all of the course material for our MAT 1070 class this semester. It is \emph{NOT} a substitute for attending class, participating in the activities, completing the online homework, and reading the textbook. It is meant to be a supplement to all of these, allowing you to go back and review material\slash{}solutions that were confusing to you when you saw them the first time. This is the blank version of the book, which contains empty versions of all class notes and activities. These will be updated before each class so that you can download or print out the notes to write on in class. The other version of the book, available on Canvas, contains the filled-in version of these notes and the answers to worksheets after they are graded.%
\end{paragraphs}%
\begin{paragraphs}{Course Notes.}{preface-2}%
These are the notes that were covered during class time. The hand-written version of the notes that was actually presented is available directly from the course Canvas page. These notes are a typed version that covers the same content in a similar level of detail.%
\end{paragraphs}%
\begin{paragraphs}{In-Class Activities.}{preface-3}%
Copies of any in-class activities are available here. This is not a substitute for working on the activity in-class, but if you would like a blank copy to reattempt later in the semester, you can find it here.%
\end{paragraphs}%
\begin{paragraphs}{Handouts.}{preface-4}%
Any helpful handouts or notes sheets will be included here.%
\end{paragraphs}%
\begin{paragraphs}{Homework.}{preface-5}%
Any hand-in homework that does not come directly from an in-class activity will be available here. The online homework is a separate part of the course, available from our Canvas site.%
\end{paragraphs}%
\end{preface}
%% begin: table of contents
%% Adjust Table of Contents
\setcounter{tocdepth}{1}
\renewcommand*\contentsname{Contents}
\tableofcontents
%% end:   table of contents
\mainmatter
%
%
\typeout{************************************************}
\typeout{Chapter  Blank Course Notes}
\typeout{************************************************}
%
\begin{chapterptx}{Chapter}{Blank Course Notes}{}{Blank Course Notes}{}{}{course-notes}
\renewcommand*{\chaptername}{Chapter}
%
%
\typeout{************************************************}
\typeout{Section  Week 1}
\typeout{************************************************}
%
\begin{sectionptx}{Section}{Week 1}{}{Week 1}{}{}{notes-week-01}
\begin{introduction}{}%
This is an outline of the topics we covered in the first week of class.%
\end{introduction}%
%
%
\typeout{************************************************}
\typeout{Handout  Monday 8\slash{}31}
\typeout{************************************************}
%
\newgeometry{left=0.75in, right=0.75in, top=0.75in, bottom=0.75in}
\begin{handout-subsection}{Handout}{Monday 8\slash{}31}{}{Monday 8\slash{}31}{}{}{notes-week-01-3}
\begin{objectives}{Objectives}{notes-week-01-3-2}
%
\begin{itemize}[label=\ptxlistdisc]
\item{}Introduction to the class%
\item{}Introduction to graphing solutions of equations.%
\end{itemize}
\end{objectives}
\begin{paragraphs}{Introduction to the Course.}{notes-week-01-3-3-1}%
Welcome to College Algebra!%
\par
What are the goals of this course?%
\par
%
\par\rule{\ptxworkspacestrutwidth}{4in}%
\end{paragraphs}%
\clearpage
\begin{paragraphs}{Introduction to the People.}{notes-week-01-3-4-1}%
Who am I?%
\par\rule{\ptxworkspacestrutwidth}{3in}%
\par
Who are all of you? Take about 3 minutes to introduce yourself to the people around you. These are going to be your companions on this semester-long journey that we have together. When we come back together, I'll want everyone to introduce themselves with their name and intended major\slash{}career path.%
\par\rule{\ptxworkspacestrutwidth}{1in}%
\end{paragraphs}%
\clearpage
\begin{paragraphs}{Syllabus and Canvas Overview.}{notes-week-01-3-5-1}%
First off, the syllabus. The entire syllabus is available on our Canvas site.%
\par
%
\begin{itemize}[label=\ptxlistdisc]
\item{}Basic info%
\item{}Qualifying%
\item{}eBook (more on that within Canvas later)%
\item{}Attendance%
\item{}Student Hours%
\item{}Homework%
\item{}Quizzes and Tests%
\item{}Grade Policy and Final Exam Requirements%
\item{}AI and Calculators%
\end{itemize}
%
\par
For the Canvas site:%
\begin{itemize}[label=\ptxlistdisc]
\item{}Home page%
\item{}Modules%
\item{}Agenda page with notes%
\item{}Pearson homework and eBook%
\end{itemize}
%
\end{paragraphs}%
\clearpage
\begin{paragraphs}{Graphing.}{notes-week-01-3-6-1}%
We're going to start this course by talking about graphing. There are many different types of graph, but the one we want to talk about in this class is the kind the allows us to compare two quantities. We want to see how the second quantity changes when the first is modified. For example:%
\par
~%
\par\rule{\ptxworkspacestrutwidth}{2in}%
\par
In order to visualize these things, we will need a two-dimensional image, where one direction corresponds to each quantity.%
\par
~%
\par\rule{\ptxworkspacestrutwidth}{2in}%
\end{paragraphs}%
\clearpage
\begin{example}{Example}{}{notes-week-01-3-7-1}%
Plot and label the points \((2,1)\), \((-4,3)\), \((0,0)\), \((-3,-2)\), \((3,-4)\), and \((2,-1)\).%
\par\rule{\ptxworkspacestrutwidth}{3in}%
\end{example}
\clearpage
\begin{paragraphs}{Solutions of Equations.}{notes-week-01-3-8-1}%
Now, what do we want to graph? Well, we're going to want to get graphs for mathematical equations (and later, functions). First, we need to look at solutions to these equations. If I have an equation in two variables \(x\) and \(y\), then a \emph{solution} is%
\par\rule{\ptxworkspacestrutwidth}{1in}%
\begin{example}{Example}{}{notes-week-01-3-8-1-3}%
Determine whether each ordered pair is a solution of the equation \(3x - 2y = 5\).%
\begin{enumerate}[font=\bfseries,label=(\alph*),ref=\alph*]%
\item{}\((1, 4)\)%
\par\rule{\ptxworkspacestrutwidth}{2in}%
\item{}\((3,2)\)%
\par\rule{\ptxworkspacestrutwidth}{2in}%
\end{enumerate}%
\end{example}
\end{paragraphs}%
\clearpage
\begin{paragraphs}{Graphs of Equations.}{notes-week-01-3-9-1}%
What is a graph of an equation?%
\par\rule{\ptxworkspacestrutwidth}{3in}%
\par
Graphs of equations that look like \(Ax+By = C\) are straight lines. Since it only takes two points to determine a straight line, we can draw these graphs by finding any two points. Two convenient points we can often use are the \emph{intercepts}.%
\begin{definition}{Definition}{}{def-intercepts}%
The intercepts of a graph are the places where the graph crosses the axes.%
\par
~%
\par\rule{\ptxworkspacestrutwidth}{3in}%
\end{definition}
\end{paragraphs}%
\clearpage
\begin{example}{Example}{}{notes-week-01-3-10-1}%
Graph \(3x - 2y = 5\).%
\par\rule{\ptxworkspacestrutwidth}{3in}%
\end{example}
\clearpage
Another way to graph equations of lines is to solve for \(y\) first, then plug points into \(x\). In the normal way we will set up these equations in the future, we'll be looking to analyze how \(y\) changes when we plug in different values of \(x\), so solving for \(y\) is usually the way to go.%
\begin{example}{Example}{}{notes-week-01-3-11-2}%
Graph \(5x + 3y = 18\).%
\par\rule{\ptxworkspacestrutwidth}{3in}%
\end{example}
\clearpage
When we solve out for \(y\) like this, we are indicating that the value of \(y\) depends on \(x\), so that \(y\) is the \emph{dependent variable} and \(x\) is the \emph{independent variable}.%
\begin{example}{Example}{}{notes-week-01-3-12-2}%
Graph \(y=x^2  - 3x + 1\).%
\par\rule{\ptxworkspacestrutwidth}{3in}%
\end{example}
Proceed to Worksheet 01 on Graphing Equations.%
\end{handout-subsection}
\restoregeometry
\end{sectionptx}
\end{chapterptx}
%
%
\typeout{************************************************}
\typeout{Chapter  Blank In-Class Activities}
\typeout{************************************************}
%
\begin{chapterptx}{Chapter}{Blank In-Class Activities}{}{Blank In-Class Activities}{}{}{activities}
\renewcommand*{\chaptername}{Chapter}
%
%
\typeout{************************************************}
\typeout{Worksheet  Graphing Equations}
\typeout{************************************************}
%
\newgeometry{left=0.75in, right=0.75in, top=0.75in, bottom=0.5in}
\begin{worksheet-section}{Worksheet}{Graphing Equations}{}{Graphing Equations}{}{}{wksht-graphing-equations}
\begin{introduction}{}%
%
\begin{itemize}[label=\ptxlistdisc]
\item{}\ptxlititle{Dates.}\par%
Assigned Date: August 31, 2026%
\par
Due Date: September 9, 2026%
\item{}\ptxlititle{Purpose.}\par%
The skill that you will practice in this exercise is graphing equations. Mastering this skill is vital to your development as a student of mathematics. Our textbook shows graphical solutions to most problems alongside the algebraic solution. Having a firm grasp of the basics will enhance a deeper understanding as you progress through the semester and on into Calculus.%
\item{}\ptxlititle{Task.}\par%
%
\begin{itemize}[label=\ptxlistcircle]
\item{}Read each question carefully.%
\item{}Answer each question fully, using complete sentences in your explanations.%
\end{itemize}
%
\item{}\ptxlititle{Resources.}\par%
Use your notes from the first class, the online eText or contact your instructor when you have questions.%
\item{}\ptxlititle{Criteria for Success.}\par%
Your assignment is due to your instructor by the date listed above. You must submit 1 pdf document as described in the syllabus, through Canvas. Your instructor will be assessing completeness and correctness for each question. This assignment will be worth 20 points.%
\end{itemize}
%
\end{introduction}%
\begin{divisionexercise}{1}{}{}{wksht-graphing-equations-3}%
[2] What does the graph of an equation represent?%
\end{divisionexercise}%
\begin{divisionexercise}{2}{}{}{wksht-graphing-equations-4}%
[9] Consider the equation \(2x+3y=12\). \emph{Fully explain} each method and your solution.%
\begin{enumerate}[font=\bfseries,label=(\alph*),ref=\alph*]%
\item{}Graph the equation using any two points.%
\item{}Graph the equation using intercepts.%
\item{}Graph the equation using the slope and \(y\)-intercept.%
\end{enumerate}%
\end{divisionexercise}%
\begin{divisionexercise}{3}{}{}{wksht-graphing-equations-5}%
[2] Which of the three ways above is your preferred method when you need to graph a \emph{linear} equation? Why?%
\end{divisionexercise}%
\begin{divisionexercise}{4}{}{}{wksht-graphing-equations-6}%
[3] Graph the equation \(y = 2x^2 - 4x + 5\).%
\end{divisionexercise}%
\begin{divisionexercise}{5}{}{}{wksht-graphing-equations-7}%
[4] Explain the method you used to graph your \emph{non-linear} equation and why another method would not be appropriate. (For example: why would graphing just two points not work?)%
\end{divisionexercise}%
\end{worksheet-section}
\restoregeometry
\end{chapterptx}
%
%
\typeout{************************************************}
\typeout{Chapter  Handouts}
\typeout{************************************************}
%
\begin{chapterptx}{Chapter}{Handouts}{}{Handouts}{}{}{handouts}
\renewcommand*{\chaptername}{Chapter}
\end{chapterptx}
%
%
\typeout{************************************************}
\typeout{Chapter  Blank Homework}
\typeout{************************************************}
%
\begin{chapterptx}{Chapter}{Blank Homework}{}{Blank Homework}{}{}{homework}
\renewcommand*{\chaptername}{Chapter}
\end{chapterptx}
\end{document}
